% ============================================================
% SECCIÓN 2: Arquitectura del Firmware
% ============================================================
\chapter{Arquitectura del Firmware}

\section{Visión General}

El firmware de Demeter implanta una red distribuida de microcontroladores ESP32 que operan
en cooperación para cubrir tres funciones complementarias: recolectar datos del entorno,
actuar sobre él y servir de puente entre el mundo físico y el sistema de información central.

Cada dispositivo ejecuta exactamente uno de tres roles, definido en tiempo de compilación:
\textbf{Nodo Sensor}, \textbf{Nodo Actuador} o \textbf{Nodo Gateway}.

% ─────────────────────────────────────────────────────────────
\section{Los Tres Tipos de Nodo}

\subsection{Nodo Sensor}

El Nodo Sensor es la unidad de observación del sistema. Su único propósito es medir el entorno
y comunicar esas lecturas al resto de la red.

\textbf{Capacidades:}
\begin{itemize}[leftmargin=*, label=--]
    \item Mide temperatura del aire, humedad relativa y humedad del suelo.
    \item Opera en modo \textbf{Activo} (reporte periódico cada N segundos) o
          \textbf{IDLE} (escucha pasiva a petición del Gateway).
    \item Soporta datos simulados para desarrollo sin hardware físico.
    \item Realiza un proceso de establecimiento de sesión (handshake) con el Gateway antes
          de iniciar el envío de telemetría.
    \item Entre reportes entra en modo \textbf{Deep Sleep}, consumiendo $\mu$A de corriente
          y alargando la vida de la batería a más de un año.
\end{itemize}

\textbf{Ciclo de vida:} El nodo se enciende, configura sus periféricos y establece contacto
con el Gateway. Una vez confirmada la sesión, alterna entre reporte y Deep Sleep.


\subsection{Nodo Actuador}

El Nodo Actuador es la unidad de acción del sistema. Recibe órdenes y las traduce en señales
físicas que controlan los elementos del invernadero.

\textbf{Capacidades:}
\begin{itemize}[leftmargin=*, label=--]
    \item Controla pines de salida digital conectados a relés (bombas, válvulas, motores).
    \item Modos de seguridad: \textbf{ACTIVO} (ejecuta comandos) e \textbf{IDLE}
          (ignora cualquier orden, cortacircuito lógico ante mantenimiento).
    \item Implementa \textbf{retroalimentación confirmada}: tras ejecutar un comando envía
          automáticamente un informe PIN\_REPORT de vuelta al servidor.
    \item Soporta ejecución de \textbf{secuencias multipaso temporizadas} sin bloquear
          el bucle principal.
\end{itemize}


\subsection{Nodo Gateway}

El Nodo Gateway es la unidad de interconexión. No mide ni actúa, sino que es el intérprete entre
la red de radiofrecuencia local y la Raspberry Pi.

\textbf{Capacidades:}
\begin{itemize}[leftmargin=*, label=--]
    \item Habla simultáneamente ESP-NOW (RF) y UART (serial), actuando como convertidor
          de medio físico sin modificar el contenido de los mensajes.
    \item Mantiene una tabla de rutas: asocia el identificador numérico de cada nodo
          con su dirección MAC física.
    \item Reenvía comandos de la Raspberry Pi (formato binario por UART) al nodo destino
          correcto en la red ESP-NOW, sin interpretarlos.
    \item Escucha todos los mensajes de la red y los reenvía transparentemente hacia la
          Raspberry Pi por UART.
\end{itemize}


% ─────────────────────────────────────────────────────────────
\section{Arquitectura Interna en Capas}

El firmware está organizado en cuatro capas claramente separadas por responsabilidad:

\begin{table}[H]
    \centering
    \caption{Capas del firmware Demeter}
    \begin{tabular}{clp{8cm}}
        \toprule
        \textbf{Nivel} & \textbf{Capa} & \textbf{Responsabilidad} \\
        \midrule
        4 (Alto) & Lógica de Nodo   & Identidad y comportamiento específico del rol. \\
        3        & Orquestador      & Máquina de estados, handshake, secuenciador. \\
        2        & Motor de Protocolo & Serialización/deserialización y enrutamiento. \\
        1 (Bajo) & Estrategia de Transporte & Canal físico: ESP-NOW, UART o mixto. \\
        \bottomrule
    \end{tabular}
    \label{tab:capas_firmware}
\end{table}

\subsection{El Orquestador del Sistema}

El orquestador es el componente más complejo del firmware. Funciona como director de orquesta
gestionando:

\begin{itemize}[leftmargin=*, label=--]
    \item \textbf{Máquina de estados:}
          \texttt{BOOT → IDLE → HANDSHAKE → RUNNING → ERROR}.
          Ninguna acción se ejecuta si el estado no lo permite.
    \item \textbf{Establecimiento de sesión} inspirado en TCP de tres pasos
          (SYN → SYN-ACK → ACK) con timeout de 2 segundos y hasta 3 reintentos.
    \item \textbf{Gestión de secuencias:} ejecución no bloqueante de listas ordenadas
          de pasos con temporización individual, evaluada en cada iteración del bucle.
    \item \textbf{Sistema de listeners (callbacks):} permite suscribirse a eventos
          sin modificar la lógica central del orquestador.
\end{itemize}

% ─────────────────────────────────────────────────────────────
\section{El Motor de Protocolo}

El Motor de Protocolo (Capa 2) actúa como el \textbf{corazón de la comunicación} del firmware.
Es completamente agnóstico al canal físico subyacente: el mismo motor gestiona tramas 
tanto si provienen de ESP-NOW como de UART.

\subsection{Responsabilidades del Motor}

\begin{enumerate}
    \item \textbf{Serialización}: transforma un comando de alto nivel (estructura de datos)
          en una trama binaria válida con cabecera, payload codificado en Little Endian y CRC.
    \item \textbf{Deserialización}: analiza una secuencia de bytes entrante, valida el
          marcador de inicio (SYNC \texttt{0xAA}), verifica el CRC y extrae el payload
          según el tipo de comando (\texttt{CMD\_ID}).
    \item \textbf{Enrutamiento}: trabaja con la tabla de rutas (ID numérico $\leftrightarrow$
          MAC física) para determinar el destinatario correcto de cada trama.
    \item \textbf{Clasificación}: identifica el tipo de mensaje recibido y lanza el
          evento interno correspondiente (callback listener).
\end{enumerate}

\subsection{Pipeline de Recepción}

Cuando el motor recibe bytes (desde cualquier transport), sigue este pipeline de validación:

\begin{enumerate}
    \item Detección del byte de sincronía (\texttt{0xAA}).
    \item Lectura de cabecera completa (5 bytes adicionales).
    \item Verificación de que el destino (\texttt{DST\_ID}) coincide con este nodo o es broadcast.
    \item Lectura del payload de longitud declarada (\texttt{LENGTH} bytes).
    \item Cálculo y comparación del CRC. Si no coincide: \textbf{trama descartada silenciosamente}.
    \item Dispatch al listener registrado para el \texttt{CMD\_ID} recibido.
\end{enumerate}


% ─────────────────────────────────────────────────────────────
\section{Estrategias de Transporte}

La Capa 1 del firmware abstrae el canal físico mediante el patrón \textbf{Estrategia}
(\textit{Strategy Pattern}): el Motor de Protocolo (Capa 2) no sabe nada sobre el
canal físico subyacente. Simplemente llama a \texttt{transport.send(bytes)} y
registra \texttt{transport.on\_receive(callback)}.

Existen tres implementaciones de transporte:

\subsection{Transporte ESP-NOW}

ESP-NOW es el protocolo propietario de Espressif para comunicación directa entre
microcontroladores ESP32 sin necesidad de punto de acceso Wi-Fi. Es la columna vertebral
de la \textbf{red de campo} de Demeter.

Características en el contexto de Demeter:
\begin{itemize}[leftmargin=*, label=--]
    \item \textbf{Alcance}: hasta 200--400 metros en campo abierto con antenas externas.
    \item \textbf{Latencia}: sub-milisegundo en condiciones normales.
    \item \textbf{MTU}: máximo 250 bytes por trama; suficiente para cualquier trama Demeter típica.
    \item \textbf{Sin confirmación}: ESP-NOW no garantiza entrega; la Capa de Aplicación aporta
          el ACK mediante el protocolo Demeter.
    \item \textbf{Direccionamiento}: por dirección MAC de 6 bytes, mapeada a un ID numérico
          en la tabla de rutas del Gateway.
\end{itemize}

Los nodos Sensor y Actuador usan ESP-NOW como canal \textbf{exclusivo}. El Nodo Gateway
usa ESP-NOW como canal de \textit{campo} y UART como canal de \textit{troncal}.

\subsection{Transporte UART (Serial)}

UART es la conexión física cableada entre el \textbf{Nodo Gateway ESP32} y la
\textbf{Raspberry Pi}, mediante los pines TX/RX a 115200 baudios.

Características en el contexto de Demeter:
\begin{itemize}[leftmargin=*, label=--]
    \item \textbf{Fiabilidad}: al ser cableado, prácticamente sin pérdida de tramas.
    \item \textbf{Velocidad}: 115200 bps → capacidad de $\approx$11.500 bytes por segundo,
          vastamente superior a la carga de telemetría esperada.
    \item \textbf{Transparencia}: el Gateway reenvía las tramas binarias íntegras sin
          interpretarlas, actuando como pasarela pura.
    \item \textbf{Gestión en ambos extremos}: el ESP32 escribe/lee sobre el UART hardware;
          la Raspberry Pi usa \texttt{pyserial-asyncio} de forma no bloqueante.
\end{itemize}

\subsection{Transporte Mixto (Gateway)}

El Nodo Gateway implementa \textbf{dos transportes simultáneos}: escucha en ESP-NOW y
retransmite los bytes recibidos por UART hacia la Raspberry, y viceversa. Esta operación
de \textit{relay} es completamente transparente: el Motor de Protocolo de la Raspberry
ve las tramas como si provinieran directamente de los nodos de campo.

\begin{designnote}
La modularidad de la capa de transporte permite añadir nuevos medios físicos (LoRa, BLE, CAN Bus)
sin modificar el Motor de Protocolo ni la Lógica de Nodo. Solo se implementa la interfaz
\texttt{TransportStrategy} y se registra en el orquestador.
\end{designnote}
